\section{SWE-bench: Benchmark cho real-world software engineering}
\label{sec:swe_bench}

\textbf{SWE-bench}~\cite{swebench} là case study tiêu biểu cho benchmark
thế hệ mới: thay vì hỏi mô hình \emph{về} lập trình, nó yêu cầu mô hình \emph{thực
sự sửa} một lỗi phần mềm có thật và kiểm tra kết quả bằng cách chạy test.

\subsection{Bài toán và cách dựng dữ liệu}

Mỗi mẫu (task instance) trong SWE-bench được trích từ một \emph{pull request} có
thật trên GitHub, thuộc các kho mã Python lớn và phổ biến (django, scikit-learn,
matplotlib, sympy\ldots). Một mẫu gồm:

\begin{itemize}
  \item \textbf{Mô tả vấn đề}: nội dung issue gốc bằng ngôn ngữ tự nhiên.
  \item \textbf{Trạng thái kho mã} tại commit ngay trước khi vá lỗi (codebase
    thật, hàng nghìn file).
  \item \textbf{Bộ test ẩn}: các unit test đi kèm PR --- gồm nhóm
    \texttt{FAIL\_TO\_PASS} (đang đỏ, phải xanh sau khi vá) và nhóm
    \texttt{PASS\_TO\_PASS} (đang xanh, không được làm hỏng).
\end{itemize}

Mô hình nhận mô tả vấn đề và kho mã, phải sinh ra một \emph{patch} (diff). Như vậy
đầu vào lấy nguyên từ thực tế, không do người ra đề tạo --- giảm khả năng benchmark
bị tối ưu hẹp và bám sát phân phối công việc kỹ sư phần mềm thật.

\subsection{Chấm bằng thực thi, không so chuỗi}

Điểm cốt lõi: SWE-bench \emph{không} so patch của mô hình với patch của con người
theo từng ký tự. Patch của mô hình được \textbf{áp vào kho mã rồi chạy lại toàn bộ
test}. Một mẫu được tính là \emph{resolved} chỉ khi: (i) mọi test
\texttt{FAIL\_TO\_PASS} chuyển sang xanh, \emph{và} (ii) mọi test
\texttt{PASS\_TO\_PASS} vẫn xanh. Tiêu chí pass/fail này khách quan và cho phép
nhiều lời giải khác nhau về mặt cú pháp nhưng cùng \emph{đúng về hành vi} --- đúng
điều mà so khớp chuỗi không làm được.

Cơ chế ``đúng = chạy được'' này chính là tinh thần ta tái hiện ở quy mô nhỏ trong
demo: ở GSM8K, ta không so chuỗi thô mà \emph{trích đáp án số rồi so giá trị}
(Phần~\ref{sec:metrics}); và ta cho thấy chính \emph{lớp trích xuất/chấm} mới là
chỗ benchmark dễ sai.

\subsection{Vì sao SWE-bench ``khó'' và có ý nghĩa}

Khi mới công bố, tỉ lệ giải được của các mô hình mạnh nhất rất thấp (cỡ vài phần
trăm), trái ngược hẳn với điểm cao chót vót trên benchmark trắc nghiệm. Sự tương
phản này minh hoạ hai luận điểm:

\begin{itemize}
  \item \textbf{Chống bão hoà}: vì đầu vào thực tế đa dạng và lời giải dài, đề
    không bị ``đụng trần'' ngay như GLUE/MMLU --- còn nhiều dư địa phân biệt.
  \item \textbf{Construct validity cao hơn}: pass/fail dựa trên test thật đo gần
    đúng \emph{năng lực sửa lỗi} hơn là một proxy trắc nghiệm.
\end{itemize}

SWE-bench cũng có các biến thể thu hẹp/được kiểm định thủ công (ví dụ tập rút gọn
và tập được con người xác minh) nhằm giảm nhiễu và chi phí --- bản thân việc này
là một bài học BetterBench~\cite{betterbench}: một benchmark tốt cần được
\emph{soi và tinh chỉnh} chứ không bất biến.

\subsection{Hạn chế và liên hệ tới demo}

SWE-bench không miễn nhiễm vấn đề: nguy cơ \emph{nhiễm dữ liệu} (mô hình từng thấy
chính PR đó khi huấn luyện), chi phí dựng môi trường thực thi lớn, và rủi ro
``đúng vì lý do sai'' (patch qua test nhưng nhờ test yếu). Đây đúng những loại rủi
ro mà demo của nhóm thu nhỏ và phơi bày: nhiễm dữ liệu (Phần~\ref{sec:results}),
độ giòn của lớp chấm (Phần~\ref{sec:metrics},~\ref{sec:error}), và sự khác biệt
giữa \emph{metric} và \emph{construct}. Nói cách khác, demo là một ``SWE-bench thu
nhỏ về mặt tư duy đo lường'': không dựng môi trường code, nhưng giữ nguyên câu hỏi
trung tâm --- \emph{phép đo có thực sự đo đúng thứ ta nghĩ không?}
